\section{Related Work \& Challenges}

\label{sec:related}

The growing use of machine learning for steady-state grid analysis has produced promising neural solvers for Power Flow (PF), Optimal Power Flow (OPF), and State Estimation (SE). Our work is motivated by four outstanding challenges: (i) unifying grid representations and model architecture across tasks, (ii) balancing physical feasibility with runtime and architectural efficiencies, (iii) improving robustness to realistic grid variations, and (iv) standardizing datasets, metrics, and workflows.

\subsection{From Task-Specific Solvers to Unified Grid Representations and Architectures}

Existing neural solvers are task-specific, which hinders the transfer of innovation; e.g., CANOS~\cite{piloto2024} pioneered the use of supervised models for large-scale OPF, but adapting it to PF required a separate redesign in \pfdelta~\cite{NEURIPS2025_d000ef56}; GNS~\cite{DONON2020106547} introduced message passing as iterative corrections for PF, but extending it to cost minimization for OPF remains open.\\

The three core steady-state tasks considered here (PF, OPF, SE) differ in formulation: PF computes bus voltages and branch flows from injections~\cite{4499318}; OPF optimizes generator setpoints given loads while considering PF equations and operational constraints~\cite{pandya2008,low2014, cain2012}; SE infers system states from sparse measurements~\cite{abur2004power}. All share electrical state variables, grid parameters, and AC physics; yet, current methods use task-specific representations that obscure this structure: PF models use bus-level graphs~\cite{powermodels, lin2024powerflownet}, OPF requires generator-level resolution~\cite{piloto2024}, and SE works with net injections rather than distinguishing generation from load~\cite{10408467}.\\

The consequences are significant. Developing neural solvers is largely an empirical process. Task-specific designs require this effort to be repeated from scratch for each problem, and advances for one task do not transfer to others. More fundamentally, no existing architecture can process the full range of inputs required across tasks and simultaneously handle cost minimization, constraint satisfaction, and robustness to noise -- the three distinct computational demands that arise across tasks. 

\subsection{Feasibility--Optimality--Runtime Tradeoffs}

A central challenge for neural grid solvers is the trade-off between runtime and solution feasibility, and in OPF, optimality. The goal of neural solvers is to reduce computation time by replacing iterative procedures (e.g., Newton–Raphson with roughly quadratic scaling in grid size~\cite{4596138}, and interior-point methods for OPF with higher-order complexity~\cite{cubic}) with a learned direct mapping from inputs to outputs. However, achieving fast inference (ideally linear in graph size) constrains neural solvers' capacity to learn this mapping, which can reduce prediction accuracy and amplify small prediction errors into degraded feasibility and optimality. \\

Existing methods address this through integration of physics and constraints into the model. For PF, GNN-based models evolve from directly regressing numerical solver outputs~\cite{lin2024powerflownet} toward physics-informed architectures, adding power-balance residuals in the loss~\cite{LOPEZGARCIA2023105567}, or predicting only subsets of variables while reconstructing the remaining states through power-flow equations~\cite{DONON2020106547}. These approaches improve physical consistency, but feasibility remains encouraged rather than guaranteed. For OPF, purely regression-based methods have also been proposed~\cite{arowolo2025}, but soft-constrained approaches have become increasingly prevalent, promoting feasibility through penalty terms in the loss function\footnote{Optimality is encoded either in a supervised way, matching pre-computed solutions~\cite{piloto2024}, or in a self-supervised one, minimizing generation costs~\cite{dc3}.}. These penalties may be combined with post-processing corrections to restore feasibility for selected constraints~\cite{piloto2024}, but are sensitive to penalty tuning and often require extensive hyperparameter search\footnote{Primal-dual~\cite{primaldual} and augmented-Lagrangian~\cite{lumina2026} methods adapt penalties dynamically but introduce more complex and potentially unstable training.}. In contrast, hard-constrained architectures enforce feasibility structurally, e.g., via projection or feasibility-seeking layers that map predictions onto the constraint manifold~\cite{NEURIPS2025_3874e2be}. While this guarantees feasibility, it also introduces higher computational costs and greater implementation complexity.\\

This feasibility--runtime trade-off is closely tied to the model architecture. Standard message-passing GNNs face a locality bottleneck: information is exchanged only between neighboring nodes at each layer, so that capturing the influence of electrically distant buses requires deeper models to increase the receptive field, which can suffer from oversmoothing and oversquashing~\cite{gnnreview, liu2022graph}, while also increasing inference cost. Model architectures with a global view, such as full-attention Transformers~\cite{vaswani2017attention, gps} reduce this locality limitation, but scale quadratically with grid size. Sparse, sampled, or hierarchical aggregation schemes improve efficiency, but introduce additional modeling choices and hyperparameters. Overall, neural grid solvers must balance feasibility, optimality, and runtime under the constraints imposed by information propagation.

\subsection{Generalization to Diverse Scenarios}

Robustness and generalization remain a central challenge for neural PF, OPF, and SE solvers. Distribution shifts in load, topology, grid parameters, and measurements can substantially degrade performance outside the training regime~\cite{NEURIPS2025_d000ef56}, unlike classical methods whose limitations stem primarily from numerical issues~\cite{tostado2021solving}.
\\

Neural solvers are largely trained on synthetic datasets with limited variation in operating conditions\footnote{Generating sufficiently large and diverse datasets is computationally expensive due to the need to repeatedly solve PF, OPF, or SE problems. For example, PGLearn required 1245 CPU Core-days for data generation~\cite{pglearn}.}. These existing datasets typically rely on uncorrelated random load perturbations around nominal operating points~\cite{lovett2024opfdatalargescaledatasetsac,varbella2024powergraph,lin2024powerflownet,LOPEZGARCIA2023105567,DONON2020106547}, fixed grid parameters~\cite{opflearn,NEURIPS2025_d000ef56,pglearn,lovett2024opfdatalargescaledatasetsac,varbella2024powergraph}, and limited topology modifications such as N-1/N-2 contingencies. All aforementioned OPF datasets keep generator cost coefficients fixed, effectively restricting learning to a single economic regime with unchanged merit order. Similarly, PF datasets, except \pfdelta, sample from OPF-feasible operating points, limiting exposure to stressed or out-of-nominal operating conditions. Constructing SE datasets is even more challenging because it additionally requires modeling sensor placement, noise, missing data, and bad-data patterns. \\

Recent datasets address some of these limitations: OPFData~\cite{lovett2024opfdatalargescaledatasetsac} and PGLearn~\cite{pglearn} were the first to scale to systems with more than 10k buses; OPF-Learn~\cite{opflearn}, PGLearn~\cite{pglearn}, and PowerGraph~\cite{varbella2024powergraph} introduced more diverse load sampling; and \pfdelta~\cite{NEURIPS2025_d000ef56} extends beyond OPF-feasible operating regions. However, these advances remain fragmented across many different data-generation libraries, which are additionally task-specific and not designed for generating the large-scale datasets required to train foundational models.\\

The importance of training data diversity has been reported in the literature; e.g., \pfdelta~\cite{NEURIPS2025_d000ef56} demonstrates improved generalization under N-1 training, yet generalization to unseen networks remains an open challenge~\cite{yang2026gridsfm,arowolo2025,DONON2020106547}. Despite this, motivations to improve dataset realism remain limited, as more challenging scenarios expose model weaknesses and often reduce reported performance. As a consequence, in the absence of standardized benchmarks, results remain difficult to compare across studies. Finally, another key limitation is the lack of validation between synthetic and real-world performance. While synthetic grid topologies used for data generation capture statistical properties of real network topologies, there is no established measure of (i) similarity between synthetic and real operating conditions, or (ii) transferability of performance from synthetic benchmarks to real grid deployment.

\subsection{Fragmented Benchmarking and Software Workflows}

Despite rapid progress, neural solver development remains fragmented. Beyond differences in data representations (bus- vs. element-level), architectures, and data-generation assumptions, methods often use incompatible input formats and features (e.g., bus-admittance matrices~\cite{varbella2024powergraph} vs. branch-level quantities~\cite{arowolo2025}). They further differ in their output representations and unit conventions, making direct comparisons difficult even for identical tasks.\\

Evaluation protocols are similarly inconsistent. Studies report heterogeneous metrics for feasibility (constraint violation magnitudes~\cite{arowolo2025, NEURIPS2025_3874e2be}, thresholded violation rates~\cite{piloto2024}) and optimality, and often rely on non-physically interpretable regression metrics (e.g., mean squared error)~\cite{varbella2024powergraph}. These are further computed with varying aggregation schemes (mean vs. max) and at different granularities (bus- vs. graph-level), making comparisons sensitive to implementation choices rather than model performance alone.\\

These issues are particularly pronounced in runtime benchmarking against classical solvers. Reported speedups depend on hardware, batch size, parallelization strategy, solver implementation, and processes included in the runtime analysis (e.g., data transfer, preprocessing, graph construction). Existing comparisons contrast highly parallelized neural solvers running on GPUs with single-threaded Python or MATLAB implementations (e.g., pandapower~\cite{pandapower} or MATPOWER~\cite{matpower}) that underutilize CPU hardware~\cite{piloto2024, NEURIPS2025_d000ef56, arowolo2025}. Without standardized evaluation protocols, runtime comparisons between neural and classical AC/DC solvers remain difficult to interpret fairly. These gaps motivate a unified software stack that standardizes data loading, graph construction, feature normalization, solver baselines, feasibility metrics, and runtime measurement protocols across PF, OPF, and SE.\\

